\subsection{Sprint 1}

Na Sprint 1, dei continuidade ao trabalho de estruturação do produto, principalmente no Painel Admin e no desdobramento técnico das funcionalidades. Segui atuando fortemente na construção e no refinamento das \textit{user stories} do admin, garantindo que os critérios estivessem mais claros para o time de desenvolvimento e alinhados com o que foi discutido com os clientes.

Também mantive minha atuação na criação e ajuste de tarefas para o gerenciamento de edificações em \textit{frontend} e \textit{backend}, além de acompanhar o detalhamento de pontos ligados a imagens e idiomas, que ainda exigiam refinamento. Essa etapa foi importante para transformar discussões mais conceituais em tarefas executáveis, com menor chance de interpretação diferente entre os times.

Nessa Sprint, comecei a perceber uma forte ausência nas nossas \textit{dailies} semanais, o que dificultou muito o andamento do projeto. As comunicações até aconteciam, mas de maneira oculta, os AGES III falavam entre si e não comunicavam decisões nem reportavam algum tipo de avanço, o que me trouxe um certo desconforto. Além disso, não obtivemos muitas respostas em relação aos AGES III sobre o que estava sendo feito, qual era o progresso da tarefa nas quais eles estavam supervisionando, ou até mesmo se existia algum tipo de bloqueio. As cobranças aconteciam e, se houvesse resposta, era sempre algo genérico do tipo "estamos trabalhando nisso", sem detalhamento do que estava sendo feito, o que me deixou bastante perdido e sem saber como ajudar ou como me organizar para contribuir melhor. Como parte disso, criei também uma planilha no \textit{Excel} para organizar os relatos da \textit{daily} semanal e verificar ausências, para assim ter um histórico e talvez alguma resposta do porquê alguma tarefa não tinha sido entregue.

Apesar disso, alguns AGES III acabaram por terminar algumas das tarefas remanescentes da Sprint, e conseguimos entregar algo viável para o cliente. Tivemos um mapa funcional, com as edificações e as informações básicas, o que foi um avanço importante para o projeto. No entanto, a AGES, por ser um ambiente acadêmico, deveria ser um espaço de aprendizado e desenvolvimento de habilidades, o que não aconteceu, já que alguns alunos da AGES I e II não tiveram interesse em terminar suas tarefas e indagar. No fim, o cliente ficou satisfeito com o que foi entregue, acabei conduzindo a apresentação para os clientes e fiquei feliz com o modo que conduzi, pois em outro projeto passado na Agência, os AGES IV tiveram dificuldades de mostrar o porquê da entrega ter tido débitos técnicos, falando até que era inadmissível. Dessa vez, eu expliquei de maneira clara e objetiva que estamos todos aprendendo e, mesmo com todas as dificuldades dessa sprint, isso serve apenas para nos fortalecer e saber para onde estamos indo.

A principal lição que levo dessa Sprint é a importância de ter uma comunicação mais transparente e frequente entre os AGES III e os AGES IV, para que todos possam entender o estado do projeto e contribuir de forma mais efetiva. A falta de \textit{feedbacks} regulares acabou gerando um ambiente de incerteza, onde não sabíamos exatamente o que estava sendo feito ou quais eram os próximos passos, o que impactou negativamente o ritmo do desenvolvimento. Combinamos de criar canais de ajuda no \textit{Discord} e de reportar avanços e bloqueios de forma mais detalhada nas \textit{dailies}, para que todos pudessem estar alinhados e contribuir melhor para o sucesso do projeto.
